\chapter*{引子\quad 历史的弹性}
\addcontentsline{toc}{chapter}{引子\quad 历史的弹性}
\markboth{引子\quad 历史的弹性}{引子\quad 历史的弹性}

1933年秋国共的第五次“围剿”和反“围剿”开始时，位于赣南、闽西的中央苏区正处于全盛时期。虽然在此之前长江流域与中央苏区可以形成掎角之势的几个大的苏区如鄂豫皖、湘鄂西已经相继被国民党军占领，但当时仍然不会有太多人想到，一年之后的秋天，这一中共控制的最大的苏维埃区域就会在国民党军强大压力下，随着主力红军的战略转移而易手。失败的结局使中共不得不走上长途跋涉之路，寻找继续生存、发展的机会，突围中的艰难也导致了中共在夺取政权前最重要的一次领导层变换。同样让人很难预想到的是，经历了如此惨痛的挫折，在对手看起来已是奄奄一息的中国共产党及其武装，很快又能重新振作，在陕北站住脚跟；而且这次失败实际上就是中共夺取全国政权前最后的一次战略性失败，从此中共的革命夺取政权之路大有直济沧海之势。

失败总不是件让人愉快的事，与中共从胜利走向胜利的历史进程的描述相比，对于中共历史上这样一次重要的失败经历，具体的研究和阐述不是很多，有许多问题我们尚不能得其详，简单的原则论述和具体的历史运行脉络也常常相差不可以道里计。而当我们重复当年更多的是基于政治考虑的结论，以“左”的错误为这次失败定性时，往往忽视了其中包含的历史的、社会的诸多因素。政治定性的高屋建瓴后面，被断送的可能是一个个正在具体影响着社会历史的细节，历史离开了细节，总让人想到博物馆那一具具人体骨骼，的确，那是人，但那真的还是人吗？

其实，翻开中共壮丽历史的长卷，在欣赏波澜壮阔的胜利画面之余，偶尔体味一下这一段别样的经历，也许可以别有一番滋味在心头。如果考虑到成败、祸福之变，谁又能说，这样的失败就完全没有意义呢？！就整个苏维埃运动而言，后人（虽然是外国人）曾有过中肯的评断：“尽管苏维埃运动遭到失败，但是政治、军事和社会经济活动的经验以及经受过组织和动员苏区居民的各种方式的尝试和失败的考验，使得中共到30年代中期成了东方各国共产党中唯一拥有实际上执政党经验的党，拥有绝无仅有的农村工作经验以及军政骨干的党。这（加上其他条件）也成为抗日战争年代里党员人数和武装力量较快增长和发展的基础。”\footnote{《共产国际、联共（布）与中国革命档案资料丛书》第7册，中央文献出版社，2002，第20～21页。}这样的说法放到第五次反“围剿”的这一时段中，也并非就没有针对性。苏维埃运动是中共革命过程中逼不得已也是不可或缺的阶段，中共革命本身就是一个从不可能到可能的创造奇迹的过程，因此所谓的超越阶段之类的说法更多只具有逻辑上的意义。作为中共首次独立领导的革命运动，苏维埃革命基本奠定了中共武装革命的思想和逻辑基础，建立了中国共产革命的第一个中央政权，通过对苏区的独立控制，显现出中共的政治理念、动员能力和控制艺术。事实上，虽然具体的权力结构和运作方式此后续有调整，但中共革命的几个重要原则诸如武装斗争、群众路线、土地革命、社会再造等，在这一时期已经牢固确立，由苏维埃革命开始，中共走上了武装夺取政权、革命建国的道路。

苏维埃革命高歌猛进的初期阶段，中共把革命的能动性发挥到了极致。1927年国共分裂时，中共几乎是白手起家开始武装反抗，仅仅数年后几十万人的武装即矗然挺立，缔造出国共合作共同北伐后的又一个传奇。中共在这其中表现出的让人感觉无穷无尽的能量，不仅当年的对手为之震惊，即连多年后的览史者，也很难不为之倾倒。不过，神话般的故事到1930年代中期暂时画下了一个逗号，再强的张力也有自己的极限，中共在多种境遇下实现的超常发展，到这时，似乎终于到了该停歇一下的时候了。从历史的大势看，1934年中共遭遇的挫折，以博古等为首的中共中央不甚成功的领导固然不能辞其咎，但这些被历史推上中心舞台的年轻人，其实本身也是历史的祭品。无论和共产革命中的前人或者后人相比，他们改变了的或者可以改变的东西实属有限，在滚滚的历史大潮面前，他们难以担当引领潮流的重任，更多时候乃是随波逐流。他们没有也不可能改变历史的航向。后人从他们身上看到的许多问题，既不一定是他们的造作，也不一定为他们所独有。对此，毛泽东曾在中共内部会议上中肯谈道：

\begin{quote}
对于任何问题应取分析态度，不要否定一切。例如对于四中全会至遵义会议时期中央的领导路线问题，应作两方面的分析：一方面，应指出那个时期中央领导机关所采取的政治策略、军事策略和干部政策在其主要方面都是错误的；另一方面，应指出当时犯错误的同志在反对蒋介石、主张土地革命和红军斗争这些基本问题上面，和我们之间是没有争论的。即在策略方面也要进行分析。\footnote{毛泽东：《学习与时局》，《毛泽东选集》第3卷，人民出版社，1991，第938～939页。}
\end{quote}

毛泽东的评判，主要是从中共内部着眼，而邓小平则从国共相争的大背景，透视过当年“围剿”与反“围剿”成败得失的幕后玄机：

\begin{quote}
如果有同志参加过十年苏维埃时期的内战，就会懂得这一点。那时不管在中央苏区，还是鄂豫皖苏区或湘鄂西苏区，都是处于敌人四面包围中作战。敌人的方针就是要扭在苏区边沿和苏区里面打，尽情地消耗我苏区的人力、物力、财力，使我们陷于枯竭，即使取得军事上若干胜利，也不能持久。\footnote{邓小平：《跃进中原的政治形势与今后的政治策略》，《邓小平文选》第1卷，人民出版社，1994，第97页。邓小平接下来还谈道：“要是按照毛主席的方针，由内线转到外线，将敌人拖出苏区之外去打就好了，那样苏区还是能够保持，红军也不致被迫长征。可惜‘左’倾机会主义者不这样做，中了蒋介石的计。”}
\end{quote}

中共两位超重量级人物的论断，客观、公正、独具慧眼，为我们提供了历史多样性认识的范本。

如果不是过分执著于结果的话，面对1933～1934年苏维埃革命的历史进程，冷静地想一想，中共在赣南、闽西这样一个狭小地区内，依靠极为有限的人力、物质资源，在国民党军志在必得、几倾全力的进攻下，竟然能够坚持一年之久，最后又从容撤退，本身也是足够令人惊叹的。何况，无论是事后诸葛的我辈，还是当年那些参与创造历史的人们，只要不是抱有“革命高潮”的狂热，对于此时国际国内背景下，红军在国民党统治中心区江南的可能命运，应该都或多或少会有不那么乐观的预判。中共和红军的成长，如毛泽东当年论述的，很大程度上就是利用着国民党内部的分化和统治力量不平衡所取得的。仔细排列一下当时各苏区的名称，诸如鄂豫皖、闽浙赣、湘鄂赣、湘赣、湘鄂西、鄂豫陕、川陕等，就可以发现一个有趣的现象，所有这些苏区都位于数省交界的边区，都利用着南京政府控制软弱的条件，力量、地域的因素在其发展中的独特作用，绝对不能低估。然而，这一切，到1930年代中期，已在悄悄发生着变化。随着地方实力派挑战的相继被击退，南京中央自身不断强化，对全国的控制力逐渐加强，中共可以利用的地方因素明显弱化，回旋空间被大大压缩。当国民党军大军压境、全力挤迫、志在必得时，成长中的中共最好的命运大概也就只能是顺利摆脱，韬光养晦，以求东山再起了。

这是一个中共成长壮大的时代，但远不是中共掌握政权的时代，超常的能量，也无法突破可以做、可能做、不能做的限界。中共在中央苏区的发展，本身在某种程度上就是毛泽东、朱德发挥自己的天才剑走偏锋（比如在军事上的天才创造、对力量的精准把握）的结果，从这一角度理解，中共的西走川陕，或许只不过是另一种形式的剑走偏锋。

所以，也许我们可以坐下来，平心静气，不抱成见，尽可能避开历史进程中现实需要带来的政治口水，更多通过当年的而不是后来的，描述性的而不是价值评判的历史资料，回首这一段曾经不那么愿意直面的历史。历史展现虽然不会像文学作品那样罗曼蒂克、激动人心，但却可能更有益于后人了解历史的本然进程，以从中汲取养分、获得智慧。实际上，每一代人都有他们自己的思考，面对着他们自己的问题，别人很难越俎代庖，因此，作为一个以国共第五次“围剿”和反“围剿”为聚焦点的研究，本书或许承担不起总结经验的责任，也未必真的能够提供若干教训，更多的只是想呈现一种面对历史的方式，即尽可能不在预设前提的背景下，去面对原初的过程。尽管，原初的历史是如此复杂，复杂得也许会让人感觉混乱，但光怪陆离既然提供给了世界，应该也就预备给了历史。

平心而论，即便自己的亲身经历我们也未必能洞察秋毫，何况那已经永远逝去的人和事，因此，原初的历史和我们的认知之间，恐怕总是会存在距离，所谓历史的弹性大概就是由此而来吧。在无限丰富的可能面前，历史研究者没有理由不谨慎和谦卑以对。当然，这绝不意味着放弃对历史本真的探求，在不确定的可能中戮力逼近那确定的唯一，是历史研究者无法逃避的宿命，否则，我们因何而存？！
